\section{风险与治理}

\subsection{风险识别}
常见的 risk 包括 budget 失控（token budget 突然突破）、verifier drift（score margin 下降）、trace drift（self-correction 不再提升 accuracy）以及 sampling collapse（diversity 下降）。每当这些风险出现时，Chen 建议立即回滚 sampling policy（温度、nucleus），并且用 “trace replay” 复现决策过程。

\begin{warningbox}{风险信号与初期应对}
1. Token budget > 110\% of SLO → reduce search width; 2. Verifier margin drop > 0.05 → rerun sampling with new prompts; 3. Trace iteration spike → throttle iteration depth and inspect trace log.
\end{warningbox}

\subsection{治理与演练}
在治理层面，需建立 inference-time runbook，包括 budget gatekeeper、trace simulator、post-mortem template。当 pipeline 异常时，可以用 runbook 迅速切换到 conservative mode（降低 token budget，限定 search depth）并把异常 log 录入 backlog，用于 future prompt iteration。

\begin{table}[H]
\centering
\caption{治理演练所需资产}
\begin{tabular}{p{0.32\textwidth}p{0.32\textwidth}p{0.32\textwidth}}
\toprule
Asset & Purpose & Trigger \\
\midrule
Budget gatekeeper dashboard & 显示 token/search/iteration 用量，支持 manual override & Budget SLO breach \\
Trace simulator & Replay candidate + verifier steps for debugging & Consistency drop or new failure mode \\
Post-mortem template & Capture root cause, fix, and regression test plan & Any production inference failure \\
\bottomrule
\end{tabular}
\end{table}

\subsection{本章小结}
有了清晰的风险识别、runbook 与治理资产，Inference-time pipeline 才能在多轮迭代中保持稳定，保障 prompt/search/trace 各环节按预算、按指标运行。

